\documentclass[literature.tex]{subfiles}

\begin{document}
\prose{Harry Potter a Kámen mudrců}{Joanne Rowlingová}
\teorieA{
epická próza, fantasy román (dobrodružný školní příběh s prvky coming-of-age)
}{
Současná dětská a young adult literatura s prvky moderního fantasy. Inspirace britskou tradicí školních příběhů, pohádek a mytologie.
}{
Přístupný, živý a imaginativní jazyk s bohatou slovní zásobou vymyšlených kouzelnických termínů (bradavický argot, jména jako Bradavice, Nebelvír). Mnoho popisů, humoru a dialogů. Spisovný jazyk s prvky hovorovosti u postav (Hagrid).
}{
Dialogy, opakování motivů (jizva, jméno „Ten, jehož jméno nesmíme vyslovit“), ironie a humor.
}{
Symbolismus (\texthl{jizva = osud a přežití}), metafora (Kámen mudrců = nesmrtelnost a chamtivost), alegorie (boj dobra proti zlu), personifikace kouzelných předmětů.
}{
Chronologické vyprávění v er-formě se vševědoucím vypravěčem. Rychlý spád, střídání napětí, humoru a dojemných momentů. Detailní popisy kouzelného světa kontrastují s šedivým mudlovským životem.
}{
U nástupiště plného lidí stála zářivě červená parní lokomotiva. Zepředu na ní bylo veliké označení Spěšný vlak do Bradavic, odjezd v 11 hodin. Harry se ohlédl a tam, kde předtím byl jízdenkový turniket, uviděl tepanou železnou bránu s nápisem Nástupiště devět a tři čtvrtě. Dokázal to! Kouř z lokomotivy se kroutil nad hlavami brebentícího davu a mezi nohama se všem pletly kočky nejrůznějších barev. Sovy rozladěně houkaly jedna na druhou tak nahlas, že je bylo slyšet i přes všechnu tu vřavu a šoupání těžkých kufrů.
}
\teorieB{
\wtem{Harry Potter}{hlavní hrdina, sirotek s jizvou ve tvaru blesku, statečný, loajální, skromný. Přechází z týraného chlapce k hrdinovi.}
\wtem{Ron Weasley}{věrný přítel z chudé kouzelnické rodiny, odvážný, vtipný, někdy nejistý.}
\wtem{Hermiona Grangerová}{chytrá, pilná, z mudlovské rodiny. Nejlepší studentka, často zachraňuje situaci znalostmi.}
\wtem{Albus Brumbál}{moudrý a laskavý ředitel Bradavic, ochránce Harryho, symbol dobra.}
\wtem{Profesor Quirrell}{učitel obrany proti černé magii, slabý a koktavý, ve skutečnosti posedlý Voldemortem.}
}{
Harry Potter žije u Dursleyových, kteří ho nenávidí. Na jedenácté narozeniny se dozví, že je čaroděj a je přijat do Bradavic. Seznámí se s Ronem a Hermionou, stane se členem Nebelvíru. Ve škole objeví, že někdo chce ukrást Kámen mudrců – artefakt umožňující nesmrtelnost. Společně s přáteli překoná řadu pastí (tříhlavý pes, ďábelské pasti atd.) a zjistí, že za tím stojí profesor Quirrell posedlý lordem Voldemortem. Harry kámen zachrání díky lásce své matky, která ho chrání před zlem.
}{
Chronologický děj s retrospektivami do Harryho dětství. Prostor: mudlovský svět (Privet Drive) × kouzelný svět (Diagon Alley, Bradavice, nástupiště 9¾). Čas: konec 20. století (Harrymu je 11 let).
}{}
\historie{}{}{}{}{}{
Joanne Rowlingová (*1965). Britská spisovatelka. Nápad na Harryho dostala v roce 1990 ve vlaku. Psala v kavárnách jako samoživitelka. První díl byl odmítnut mnoha nakladateli, pak uspěl u Bloomsbury. Série se stala celosvětovým fenoménem.
}{
Celá série \textsc{Harry Potter} (7 dílů), \textsc{Fantastická zvířata a kde je najít}, detektivní série pod pseudonymem Robert Galbraith (\textsc{Cormoran Strike}).
}
\kritika{}{}{
Témata přátelství, odvahy, boje se zlem a nalezení vlastní identity zůstávají nadčasová. Kniha popularizovala fantasy pro děti a mládež a inspirovala celou generaci čtenářů.
}
\nazor{
První díl Harryho Pottera je dokonalý start do kouzelného světa – chytne od první stránky. Nejvíce mě baví kontrast mezi šedivým životem u Dursleyových a barevnými Bradavicemi, stejně jako trio Harry–Ron–Hermiona, které funguje skvěle. Kámen mudrců je chytrý MacGuffin a finále s Quirrellem/Voldemortem napínavé. Rowlingová mistrně mísí humor, napětí a dojemné momenty. I po letech je to jedna z nejpřístupnějších a nejkouzelnějších fantasy sérií. Rozhodně doporučuji – otevírá bránu do světa, kde je přátelství a odvaha silnější než temná moc.
}
\explain{
\textbl{Fantasy --} žánr s prvky magie, mytologie a dobrodružství;
\textbl{Coming-of-age --} příběh dospívání a hledání identity;
\textbl{MacGuffin --} důležitý předmět, který pohání děj (Kámen mudrců).
}
\wpage
\end{document}